\section{Pseudocode}
\label{sec:pseudocode}

The \href{https://ctan.org/pkg/algorithmicx}{algorithmicx} package provides a
structured way to typeset pseudocode. For example, \cref{alg:euclidean}
computes the greatest common divisor of two non-negative integers using
Euclid's algorithm.

\begin{algorithm}[htb]
	\caption{Euclidean algorithm}
	\label{alg:euclidean}
	\begin{algorithmic}[1]
		\Require Integers $a \geq 0$ and $b \geq 0$, not both zero
		\Ensure $\gcd(a,b)$
		\While{$b \neq 0$}
			\State $r \gets a \bmod b$
			\State $a \gets b$
			\State $b \gets r$
		\EndWhile
		\State \Return $a$
	\end{algorithmic}
\end{algorithm}

Use \verb|\caption{...}| and \verb|\label{...}| as for other floats, and
refer to the algorithm with \verb|\cref{...}|. The optional argument
\verb|[1]| on the \texttt{algorithmic} environment enables line numbers.
